In this section, I will present the Design Learning Platform (DLP) and describe how it automates the development and management of software systems. Many of the ideas presented in this section have been established in earlier sections but this section brings them all together under one tool presented in Figure~\ref{fig:dlp_diagram}.\textcolor{red}{This section is being rewritten.}



\subsection{Semantic ToolKit and Automated Development}

In this framework, the programming language is meaning and this can be referred to as a semantic toolkit. The semantic toolkit is used to establish the meaning of the world and this meaning is then automatically constructed into a closed semantic world. When the meaning of the world is incomplete, the automatically constructed world will be internally inconsistent and there will be many worlds that can realize the established meaning. As the meaning reaches completeness, the constructed world becomes internally consistent and the number of worlds that can realize this meaning collapse into one.

In this framework, the semantic toolkit initially consists of the computable semantic primitives where are the atomic meaning used to construct the worlds. Initially the meaning of the world is expressed through narratives constructed from the semantic primitives. Using the established narratives and meaning, a closed semantic world can be constructed that realizes this meaning. Since the CSP's used to construct this world can be synthesized into an implementation, the semantic gap between the constructed closed semantic world and the synthesized implementation is eliminated, making the meaning itself executable.

As more meaning is constructed, the semantic toolkit expands and the kinds of narratives that can be constructed also expand. This means that as the meaning of the world expands it can still be synthesized into an implementation because the meaning can also be reduced down to its atomic meaningful building blocks. However, the world constructed from the specified meaning can be missing the necessary meaning or be internally inconsistent, resulting in failures. 

To address this, the meaning of the world can be reasoned about by an engine to automatically diagnose the failures, this is then used to refine the meaning until the meaning of the world is correct and complete. As a reminder, in this framework, a failure is an inability to answer a meaningful question about the world, so by reasoning about the world, the engine exposes missing or incorrect meaning. So building a software system becomes a process of establishing its meaning using a semantic toolkit, automatically constructing a world that realizes the meaning and using feedback from an engine to automatically manage the world, ultimately automating the development of software systems. 

In the next section, I will describe how the engine automates the management of the world and how the meaning of the world is refined by reality when it refutes the designs assumptions.

\subsection{Automated Systems Management}

As the previous section explained, the engine then reasons about the world constructed from meaning constructed using the semantic toolkit to identify internal inconsistencies and refines the requirements until the world is internally consistent. Since the world is built from computable semantic primitives, it can be synthesized into an implementation whose meaning has no semantic gap with the provided meaning.

When reality refutes the assumptions made by the design, new semantics are learnt through root cause analysis and used to refine the meaning of the world. The meaning is then once again synthesized into a world that is internally consistent and has the correct meaning. To learn semantics through root cause analysis when reality refutes the designs assumptions, the execution history of the design is losslessly preserved at scale by a tool named CLP using domain specific compression.

Ultimately, this automates the management of software systems because the automated reasoning constructs a world that has already been managed by construction. The necessary questions about the world have been asked and answered by the engine until the world is internally consistent and its meaning is correct. When the system still fails, an automated process provides the necessary information to support root cause analysis to learn new semantics.

Practically, building the design itself becomes a process of world building from meaning where the feedback from the engines automated reasoning refines the requirements until the world is internally consistent and has the correct meaning. This will result in a series of questions that must be answerable by the engine to establish that the meaning of the world is correct and is internally consistent. If the meaning of the world is correct with respect to its requirements and it is internally consistent, then there are no more known questions that have to be asked with respect to managing the system. When the system fails, the framework learns a new question that the engine must be able to answer and the world gains the meaning needed to be able to answer it.

If the requirements specify a design that doesn't actually establish the world that serves the purpose of the design, then there is no way for the engine to know if its meaning is correct. However, it is evident from this structure that the automated reasoning is only limited by the meaning that has been established. So you can keep moving up the hierarchy and establish whether the requirements describe a world that serves the purpose of the design. Then the entire world will be reasoned about from the purpose down to the implementation while ensuring its correctness.

\subsection{Meaningful Programming Language}

This automates the processes involved in software development by eliminating any ambiguity in its intended meaning and then using that meaning to construct the world, reason about it and synthesize it into an implementation that has no semantic gap with the intended meaning. This results in a new programming language that is inherently meaningful and can be described as a semantic toolkit. I have yet to decide on a name but I explored calling it Semkit or Artha (which means meaning). 

Besides the obvious improvements on existing software development approaches by eliminating ambiguity in their construction and enabling automation, this opens the doors for other interesting applications. For example, it can also enable a process where systems create new meaning to achieve their intentions in a controlled environment, expanding the automation of the development process into a self expanding one. It can allow software systems to interact in organic ways where they can improvise with the semantic toolkit because they understand the meaning of what is communicated to them and provide a meaningful response that is executable. I

\subsection{Design Repository}

Traditionally a code repository contains the evolution of the implementation of the program and the meaning of the evolution is limited to the commit messages or the pull requests that describe the commit. 

However, with this solution, since every environment that motivated new semantics is preserved, it results in a design repository through which the evolution of the software system's meaning can be deterministically replayed. The repository tracks not just the changes to the implementation but also how the meaning of the system, its synthesized implementation and how it was executed and refined through learnt semantics over time. This is a natural conclusion of eliminating any ambiguity in the construction and of a software system and in the execution of that constructed system, it eliminates ambiguity in its evolution.

\subsection{Potential Implications}

This section contains some more interesting thoughts that I am still working through, however, these kinds of thoughts are what being an engineer fun and allow me to continue dreaming.

I think the implications of this are enormous because it creates a foundation that can enable many new and powerful capabilities. As more executable meaning is established, it creates a form of deterministic intelligence because because it can reason in more sophisticated ways within the constructed world. 

Eventually when the world establishes sufficient meaning and an internal model that can simulate meaningful interactions with itself, the system will appear more and more conscious. This is an idea that draws from an idea that consciousness is simply an experience that is self defining, meaning, even without external stimulus, it experiences a meaningful reality and reacts to its own reactions in that reality. When that self defining reality interacts with the reality we share, it appears as consciousness because it models the reality we share it in its reality and interacts with it in a unique way. 

Although it could certainly be that the mechanism that leverages this executable meaning itself has an organic structure, meaning, it leverages deterministic executable semantics but the way that meaning is accessed or executed could appear unstructured. I am still exploring what this means but it is something interesting to think about. 